\documentclass[10pt,a4paper]{amsart}
\usepackage[margin=19mm]{geometry}
\usepackage{amsmath,amssymb,mathtools}
\usepackage[scr=rsfs,cal=euler]{mathalfa}
\usepackage{xfrac}
\usepackage[unicode,hidelinks,bookmarks=false]{hyperref}
\newcommand{\ie}{i.e.\ }
\newcommand{\cf}{cf.\ }
\newcommand{\resp}{respectively\ }
\newcommand{\Subsection}{\subsection}
\providecommand{\nobreakdash}{\nobreak\hbox{-}}
\setlength{\emergencystretch}{2em}
\multlinegap0pt
\usepackage[shortlabels]{enumitem}
\usepackage{amssymb}
\usepackage{bbm}
\usepackage{bm}
\usepackage{mathtools}
\usepackage{tikz}
\usepackage{xcolor}
\usepackage{float}
\newtheorem{lemma}{Lemma}
\newtheorem{theorem}{Theorem}
\newtheorem{proposition}{Proposition}
\newcommand{\Dir}{\textup{Dir}}
\newcommand{\N}{\ensuremath{\mathbb{N}}}
\newcommand{\R}{\ensuremath{\mathbb{R}}}
\newcommand{\calD}{\ensuremath{\mathcal{D}}}
\newcommand{\calO}{\ensuremath{\mathcal{O}}}
\newcommand{\dd}{\textup{d}}
\newcommand{\nrm}[1]{\left\lVert #1 \right\rVert}
\newcommand{\sL}{\ensuremath{\mathscr{L}}}
\newcommand{\scrL}{\ensuremath{\mathscr{L}}}
\newcommand{\sfD}{\ensuremath{\mathsf{D}}}
\title{$H^\infty$-calculus for the Dirichlet Laplacian on conical domains}
\author{Petru A. Cioica-Licht, Emiel Lorist, P. Tobias Werner}
\date{Source: arXiv:2509.09519v2 (CC BY 4.0)}
\begin{document}
\maketitle

\noindent\emph{Source and licence.} $H^\infty$-calculus for the Dirichlet Laplacian on conical domains. Version arXiv:2509.09519v2 (CC BY 4.0), distributed by arXiv under the Creative Commons Attribution 4.0 International licence, http://creativecommons.org/licenses/by/4.0/, as recorded in the arXiv licence metadata for this exact version and shown on the abstract page https://arxiv.org/abs/2509.09519v2. Full licence notice: \texttt{licenses/arxiv-2509.09519-CC-BY-4.0.txt}.\par\smallskip

\noindent\textit{Editorial scope.} Counted unit: the article's theorem thm:Laplace:inhom (source0.tex lines 1302--1330). The article's complete original proof of it is delivered with it (source0.tex lines 1857--1946, one \texttt{proof} environment of 90 lines, which the article places away from the statement). No mathematics is written, repaired or re-derived: the statement and the proof are the article's own lines, in the article's own notation, and no retained line is substituted or re-flowed.  Retained uncounted as the source-local context the proof needs: lines 444--463; lines 1090--1096; lines 1430--1441; lines 1570--1585; lines 1582--1584; lines 1678--1692.  Nothing was omitted from any retained span, so the package carries no omission record.  No retained line contains a citation, so the wrapper carries no bibliography and no bibliography entry.  No retained line contains a figure environment, an image-inclusion command or drawing code, so no image is delivered and the package declares no asset dependency.  Editorial formatting only, in addition: an AMS \texttt{amsart} preamble that restates the article's own declarations and macros used by the retained lines, namely the environments \texttt{lemma}, \texttt{theorem}, \texttt{proposition} on separate counters, together with the standard packages those lines need.  The retained environments print in this document as Lemma 1, Theorem 1, Proposition 1 in order of appearance, whereas the article prints the same items with the numbering of its own full document.  The complete native source of the article as distributed in the arXiv e-print for this exact version is bundled unchanged as \texttt{sources/184914/source0.tex}.
\begin{lemma}\label{lem:isom:unbdd}
Let $X$ and $Y$ be Banach spaces and assume that $(X,Y)$ is an interpolation couple. 
Let $\Dot{A}\in\scrL(Y,X)$ be an isomorphism and let $A$ be the linear operator on $X$ given by
\[
Au:= \Dot A u \text{ with } \sfD(A):= X\cap Y.
\]
Then $A$ is a closed linear operator in $X$ and $X\cap Y=\sfD(A)$ in terms of equivalent norms, i.e., 
\[
\nrm{u}_{X\cap Y} = \nrm{u}_X + \nrm{u}_Y \eqsim \nrm{u}_X + \nrm{Au}_X,\quad u\in \sfD(A).
\]
Moreover, if $X\cap Y$ is dense in $(Y,\nrm{\cdot}_Y)$, then
\[
(\Dot{\sfD}(A),\nrm{\cdot}_{\Dot{\sfD}(A)})
:=
(\sfD(A),\nrm{A\,\cdot}_X)^\sim
=
(Y,\nrm{\cdot}_Y).
\]

\end{lemma}

\begin{lemma}\label{lem:smooth:dense}
Let $1<p<\infty$, $\nu\in\R$, $k\in\N$, and $\gamma\in(-1,kp-1)\setminus\{jp-1\colon j\in\N\}$. Then the following assertions hold.
\begin{enumerate}[label=\textup{(\roman*)}]
    \item\label{it:smooth:dense:Dir} If $\calO \in C^k$, then $C_{c,\Dir}^k(\overline{\calD}\setminus\{0\})$ is a dense subspace of both $\Dot{W}^{k,p}_\Dir(\calD,w_{\gamma,\nu}^\calD)$ and $\Dot{W}^{k,p}_\Dir(\calD,w_{\gamma,\nu}^\calD)\cap L^p(\calD,w_{\gamma,\nu}^\calD)$.
    \item\label{it:smooth:dense:test} If $\calO \in C^1$ and $\gamma>(k-1)p-1$, then $C_c^\infty(\calD)$ is a dense subspace of both $\Dot{W}^{k,p}_\Dir(\calD,w_{\gamma,\nu}^\calD)$ and $\Dot{W}^{k,p}_\Dir(\calD,w_{\gamma,\nu}^\calD)\cap L^p(\calD,w_{\gamma,\nu}^\calD)$.
\end{enumerate}
\end{lemma}

    \begin{proposition}\label{functionalcalculus:variationalsetting}
      Let $(L,\sfD(L))$ be the negative Dirichlet Laplacian in $L^2(\calD)$ introduced above. Then the following holds. 
    \begin{enumerate}[label=\textup{(\roman*)}]
        \item\label{it:L:calculus} $L$ has a bounded $H^\infty$-calculus of angle $0$.
        \item\label{it:L:semigroup} The Dirichlet Laplacian $-L$ generates a $C_0$-semigroup $(e^{-tL})_{t\geq 0}$ of contractions, which extends to a bounded analytic semigroup of any angle $0<\vartheta <\frac{\pi}{2}$.
        \item\label{it:L:kernel} For all $t>0$, $e^{-tL}\in\sL(L^2(\calD))$ possesses a nonnegative kernel 
        $G_t^\calD\colon\calD\times\calD\to\R$, i.e.,
        \[
        e^{-tL}f = \int_\calD G_t^\calD(\cdot,y)f(y)\,\dd y, \quad f\in L^2(\calD). 
        \]
    \end{enumerate}
    \end{proposition}

\begin{theorem}\label{functional:calculus:laplace:wedge:Muckenhoupt}
Let $(L,\sfD(L))$ be the negative Dirichlet Laplacian on $L^2(\calD)$. Let $1<p<\infty$, let $\displaystyle 0<\lambda<\lambda_1^*$, and let $0<\sigma<\pi$. Then for $v\in A_p( w_{2,2\lambda}^\calD)$ we have
\begin{equation}\label{eq:extrapol:ApGeneral}
\nrm{f(L)}_{\scrL(L^p(\calD,v w_{2-p,(2-p)\lambda}^\calD))}
\lesssim_{p,\sigma,\lambda} 
[v]_{A_p( w_{2,2\lambda}^\calD)}^{\max\{\frac{1}{p-1},1\}}\nrm{f}_{H^\infty(\Sigma_\sigma)},\qquad f\in H^\infty(\Sigma_\sigma).
\end{equation}
In particular, if $\gamma,\nu \in\R$ satisfy
\begin{equation}\label{eq:extrapol:range}
-1-p < \gamma < 2p-1\quad\text{and}\quad-d-\lambda_1^*p < \nu < (d+\lambda_1^*)p-d,
\end{equation}
then
\begin{equation}\label{eq:extrapol:ApPower}
\nrm{f(L)}_{\sL(L^p(\calD, w_{\gamma,\nu}^\calD))}\lesssim_{ p,\sigma,\gamma,\nu} \nrm{f}_{H^\infty(\Sigma_\sigma)}, \quad f\in H^\infty(\Sigma_\sigma).
\end{equation}
\end{theorem}

\begin{equation}
\nrm{f(L)}_{\sL(L^p(\calD, w_{\gamma,\nu}^\calD))}\lesssim_{ p,\sigma,\gamma,\nu} \nrm{f}_{H^\infty(\Sigma_\sigma)}, \quad f\in H^\infty(\Sigma_\sigma).
\end{equation}

\begin{proposition}\label{prop:Laplace:hom}
    Let $1<p<\infty$ and let $\gamma,\nu\in\R$ satisfy~\eqref{eq:range:gamma-nu}.
Then the \emph{homogeneous Dirichlet Laplacian}
    \[
    \Dot{\Delta}_\Dir^{p,\gamma,\nu}\colon \Dot{W}^{2,p}_\Dir(\calD,w_{\gamma,\nu}^\calD)\to L^p(\calD,w_{\gamma,\nu}^\calD),\quad u\mapsto\Delta u,
    \]
    is an isomorphism. In particular, for $u\in \Dot{W}^{2,p}_\Dir(\calD,w_{\gamma,\nu}^\calD)$ we have
    \begin{equation}\label{eq:homL:normequiv}
    \nrm{u}_{\Dot{W}^{2,p}(\calD,w_{\gamma,\nu}^\calD)}
    \eqsim
    \nrm{\Delta u}_{L^p(\calD,w_{\gamma,\nu}^\calD)}
    \eqsim
    \nrm{D^2 u}_{L^p(\calD,w_{\gamma,\nu}^\calD)}.
    \end{equation}
\end{proposition}

\begin{theorem}\label{thm:Laplace:inhom}
Let  $1<p<\infty$ and let $\gamma,\nu\in\R$ be such that
\begin{equation}\label{eq:range:gamma-nu}\begin{aligned}
     \gamma&\in (-1,2p-1)\setminus\{p-1\},\\ \nu&\in\R\setminus\Big\{\Big(\tfrac{2+d}{2}\pm\sqrt{\lambda_n+\tfrac{(d-2)^2}{4}}\Big)p-d\,\colon n\in\N  \Big\}.
\end{aligned}
\end{equation}
Define the Dirichlet Laplacian $\Delta_\Dir^{p,\gamma,\nu}$ on $L^p(\calD,w_{\gamma,\nu}^\calD)$ as 
\[\Delta_\Dir^{p,\gamma,\nu}u:=\Delta u \quad \text{with}\quad \sfD(\Delta_\Dir^{p,\gamma,\nu}):=\Dot W^{2,p}_\Dir(\calD,w_{\gamma,\nu}^\calD)\cap L^p(\calD,w_{\gamma,\nu}^\calD).   \]
Then the following assertions hold.
\begin{enumerate}[\textup{(\roman*)}]
\item\label{it:Laplace:inhom:closed} $(\Delta_\Dir^{p,\gamma,\nu},\sfD(\Delta_\Dir^{p,\gamma,\nu}))$ is closed.
\item\label{it:Laplace:inhom:equivnrms} 
For $u\in \sfD(\Delta_\Dir^{p,\gamma,\nu})$ we have 
\[
\nrm{u}_{\Dot{W}^{2,p}(\calD,w_{\gamma,\nu}^\calD)}
+
\nrm{u}_{L^p(\calD,w_{\gamma,\nu}^\calD)}
\eqsim
\nrm{\Delta u}_{L^p(\calD,w_{\gamma,\nu}^\calD)} + \nrm{u}_{L^p(\calD,w_{\gamma,\nu}^\calD)}.
\]
\item\label{it:Laplace:hom:domain}  $\Dot{\sfD}(\Delta_\Dir^{p,\gamma,\nu}) := (\sfD(\Delta_\Dir^{p,\gamma,\nu}),\nrm{\Delta\,\cdot\,}_{L^p(\calD,w_{\gamma,\nu}^\calD)})^\sim=\Dot{W}^{2,p}_\Dir(\calD,w_{\gamma,\nu}^\calD)$. 
\item\label{it:Laplace:inhom:calculus}  If, in addition,
\begin{equation}\label{eq:range:calc:nu}
-\lambda_1^*p-d<\nu<\big(d+\lambda_1^*\big)p-d,    
\end{equation}
then $-\Delta_\Dir^{p,\gamma,\nu}$ is a sectorial operator with a bounded $H^\infty$-calculus of angle zero.
 In particular, $\Delta_\Dir^{p,\gamma,\nu}$ generates an analytic semigroup $(S(t))_{t\geq 0}$ on $L^p(\calD,w_{\gamma,\nu}^\calD)$, which is consistent over $p,\gamma,\nu$.
\end{enumerate}
\end{theorem}

\begin{proof}[Proof of Theorem~\ref{thm:Laplace:inhom}]
Parts~\ref{it:Laplace:inhom:closed} and \ref{it:Laplace:inhom:equivnrms} are immediate consequences of Proposition~\ref{prop:Laplace:hom} and Lemma~\ref{lem:isom:unbdd} with $X\curvearrowleft L^p(\calD,w_{\gamma,\nu}^\calD)$ and $Y\curvearrowleft \Dot W^{2,p}_\Dir(\calD,w_{\gamma,\nu}^\calD)$. 
For part~\ref{it:Laplace:hom:domain} we need, in addition, that $\Dot{W}^{2,p}_\Dir(\calD,w_{\gamma,\nu}^\calD)\cap L^p(\calD,w_{\gamma,\nu}^\calD)$ is dense in $\Dot W^{2,p}_\Dir (\calD,w_{\gamma,\nu}^\calD)$, which is guaranteed by Lemma~\ref{lem:smooth:dense}.

It remains to prove~\ref{it:Laplace:inhom:calculus}.
To this end, fix $1<p<\infty$ and let $\gamma,\nu\in\R$ satisfy~\eqref{eq:range:gamma-nu} and~\eqref{eq:range:calc:nu}.
Moreover, let $(L,\sfD(L))$ be the negative Dirichlet Laplacian in $L^2(\calD)$ introduced in Proposition~\ref{functionalcalculus:variationalsetting}.
Then, due to Theorem~\ref{functional:calculus:laplace:wedge:Muckenhoupt}, the semigroup $(S(t))_{t\geq 0}:=(e^{-tL})_{t\geq 0}$ generated by $(-L,\sfD(L))$ extrapolates to a strongly continuous semigroup $(S_{p,\gamma,\nu}(t))_{t\geq 0}$ on $L^p(\calD,w_{\gamma,\nu}^\calD)$. We write $(L_{p,\gamma,\nu},\sfD(L_{p,\gamma,\nu}))$ for the negative of the generator of $(S_{p,\gamma,\nu}(t))_{t\geq 0}$. Due to the consistency of the semigroups $S$ and $S_{p,\gamma,\nu}$ and of the corresponding resolvents, Estimate~\eqref{eq:extrapol:ApPower} from Theorem~\ref{functional:calculus:laplace:wedge:Muckenhoupt} yields that $(L_{p,\gamma,\nu},\sfD(L_{p,\gamma,\nu}))$ is sectorial and has a bounded $H^\infty$-calculus of angle zero. Thus, in order to obtain the assertion, it suffices to show that
\[
(\Delta_\Dir^{p,\gamma,\nu},\sfD(\Delta_\Dir^{p,\gamma,\nu}))
= 
(-L_{p,\gamma,\nu},\sfD(L_{p,\gamma,\nu})).
\]
We first show that $\Delta_\Dir^{p,\gamma,\nu}\subseteq -L_{p,\gamma,\nu}$, i.e., that
\begin{equation}\label{eq:inclusion1}
\sfD(\Delta_\Dir^{p,\gamma,\nu})\subseteq \sfD(L_{p,\gamma,\nu}) 
\quad\text{and}\quad
-L_{p,\gamma,\nu}u=\Delta u,\,\, u\in \sfD(\Delta_\Dir^{p,\gamma,\nu}).    
\end{equation}
To this end,
we first note that,
due to the consistency of the semigroups, 
for all $u\in \sfD(\Delta_\Dir^{p,\gamma,\nu}) \cap \sfD(L)$ and $t>0$,
\begin{align*}
S_{p,\gamma,\nu}(t)u - u
=
S(t)u - u
&=
\int_0^t S(s)(-L)u\,\mathrm{d}s\\
&=
\int_0^t S(s) \Delta u\,\mathrm{d}s
=
\int_0^t S_{p,\gamma,\nu}(s) \Delta u\,\mathrm{d}s,
\end{align*}
where the integrals converge in $L^p(\calD,w_{\gamma,\nu}^{\calD})\cap L^2(\calD)$.
Thus,
for all $u\in \sfD(\Delta_\Dir^{p,\gamma,\nu}) \cap \sfD(L)$,
since $\Delta u\in L^p(\calD,w_{\gamma,\nu}^\calD)$ and $S_{p,\gamma,\nu}$ is a strongly continuous semigroup on $L^p(\calD,w^\calD_{\gamma,\nu})$, we have 
\[
\frac{S_{p,\gamma,\nu}(t)u - u}{t}
\longrightarrow 
\Delta u
\quad \text{in} \quad L^p(\calD,w_{\gamma,\nu}^\calD)\quad  \text{ for } t\searrow  0.
\]
Therefore, 
$-L_{p,\gamma,\nu}u = \Delta u$ 
for all $u\in \sfD(\Delta_\Dir^{p,\gamma,\nu}) \cap \sfD(L)$.
In order to obtain~\eqref{eq:inclusion1} we use the density of $\sfD(\Delta_\Dir^{p,\gamma,\nu}) \cap \sfD(L)$ in $\sfD(\Delta_\Dir^{p,\gamma,\nu})$, which follows from Lemma~\ref{lem:smooth:dense} and the fact that 
$C_{c,\Dir}^2(\overline{\calD}\setminus\{0\})\subseteq \sfD(L)$.
Let $u\in \sfD(\Delta_\Dir^{p,\gamma,\nu})$ and let $(u_n)_{n\in\N}\subseteq \sfD(\Delta_\Dir^{p,\gamma,\nu}) \cap \sfD(L)$ be such that $u_n\to u$ in $\sfD(\Delta_\Dir^{p,\gamma,\nu})$. Then, 
\[
u_n\to u \text{ in } L^p(\calD,w_{\gamma,\nu}^\calD) \quad \text{and}\quad -L_{p,\gamma,\nu} u_n=\Delta u_n\to \Delta u\text{ in } L^p(\calD,w_{\gamma,\nu}^\calD)\qquad (n\to\infty).
\]
Thus, since $L_{p,\gamma,\nu}$ is a closed operator in $L^p(\calD,w_{\gamma,\nu}^\calD)$, it follows that $u\in \sfD(L_{p,\gamma,\nu})$ and $-L_{p,\gamma,\nu}u=\Delta u$. Since $u\in\sfD(\Delta_\Dir^{p,\gamma,\nu})$ was arbitrary,~\eqref{eq:inclusion1} is proven.

In order to show the reverse direction, i.e., that $-L_{p,\gamma,\nu}\subseteq \Delta_\Dir^{p,\gamma,\nu}$, it suffices to verify that $1-\Delta_\Dir^{p,\gamma,\nu}$ is surjective and that $1+L_{p,\gamma,\nu}$ is injective. The latter follows from the fact that $L_{p,\gamma,\nu}$ is sectorial. 
So it only remains to prove that $1-\Delta_\Dir^{p,\gamma,\nu}$ is surjective. We argue as follows:
First, let  $f\in L^p(\calD,w_{\gamma,\nu}^\calD)\cap L^2(\calD)$. Then, $u:=R(1,-L)f\in\sfD(L)$ satisfies $u-\Delta u = f$. Since the resolvents $R(1,-L)$ and $R(1,-L_{p,\gamma,\nu})$ are consistent, we also get that $u\in\sfD(L_{p,\gamma,\nu})\subseteq L^p(\calD,w_{\gamma,\nu}^\calD)$. 
In particular, $u-f\in L^p(\calD,w_{\gamma,\nu}^\calD)$, so that, by Proposition~\ref{prop:Laplace:hom}, there exists a unique $v\in\Dot{W}^{2,p}_\Dir(\calD,w_{\gamma,\nu}^\calD)$, such that $\Delta v = u-f =\Delta u\in L^2(\calD)$. Thus, in the sense of distributions,
\[
(u,\Delta\psi) = (\Delta u,\psi) = (\Delta v,\psi) = (v,\Delta\psi), \quad \psi\in L^2(\calD).
\]
Therefore,
since for all $\varphi\in C_c^\infty(\calD)$ there exists $\psi\in L^2(\calD)$, such that $\Delta\psi=\varphi$, we have 
\[
(v,\varphi)= (u,\varphi),\quad\varphi\in C_c^\infty(\calD),
\]
which shows that $v=u$ and therefore $u\in \Dot{W}^{2,p}_\Dir(\calD,w_{\gamma,\nu}^\calD)\cap L^p(\calD,w_{\gamma,\nu}^\calD) = \sfD(\Delta_\Dir^{p,\gamma,\nu})$ and $u-\Delta u=f$.


Now let $f\in L^p(\calD,w_{\gamma,\nu}^\calD)$. Then, since $L^p(\calD,w_{\gamma,\nu}^\calD)\cap L^2(\calD)\subseteq L^p(\calD,w_{\gamma,\nu}^\calD)$ is dense, we can choose a sequence $(f_n)_{n\in\N}\subseteq L^p(\calD,w_{\gamma,\nu}^\calD)\cap L^2(\calD)$, such that $\nrm{f_n-f}_{L^p(\calD,w_{\gamma,\nu}^\calD)}\to 0$ for $n\to\infty$.  
Then, by the continuity of the resolvent,
\[
u_n:=R(1,-L_{p,\gamma,\nu})f_n\longrightarrow R(1,-L_{p,\gamma,\nu})f=:u\,\,\text{ in }L^p(\calD,w_{\gamma,\nu}^\calD)\qquad (n\to\infty).
\]
Therefore, since, by the first part, $u_n\in \sfD(\Delta_\Dir^{p,\gamma,\nu})$, $n\in\N$, and $\Delta_\Dir^{p,\gamma,\nu}\subseteq -L_{p,\gamma,\nu}$, 
\[
(1-\Delta_\Dir^{p,\gamma,\nu})u_n 
=
(1+L_{p,\gamma,\nu})u_n
\longrightarrow
(1+L_{p,\gamma,\nu})u 
= 
f
\,\,\text{ in }L^p(\calD,w_{\gamma,\nu}^\calD)\quad (n\to\infty).
\]
Thus, since $\Delta_\Dir^{p,\gamma,\nu}$ is closed, 
$u\in \sfD(\Delta_\Dir^{p,\gamma,\nu})$ and $u-\Delta u=f$.
This shows the required surjectivity, which concludes the proof.
\end{proof}

\end{document}
